% !TEX root = exercices_s6.tex
% ============================================================================
% HEC Montréal MBA -- ECON50803: Macroeconomic Environment
% Shared Preamble for Exercise Handouts
% ============================================================================

% --- Document class ---------------------------------------------------------
\documentclass[11pt, letterpaper]{article}
\usepackage[margin=2.5cm]{geometry}

% --- Fonts ------------------------------------------------------------------
\usepackage[sfdefault,light]{FiraSans}
\usepackage{FiraMono}
\usepackage[T1]{fontenc}

% --- Colors (same as slides) ------------------------------------------------
\usepackage{xcolor}
\definecolor{HECnavy}{HTML}{002855}
\definecolor{HECgreen}{HTML}{26d07c}
\definecolor{HECcoral}{HTML}{ff585d}
\definecolor{HECtext}{HTML}{2D3748}
\definecolor{HECgray}{HTML}{94A3B8}
\definecolor{HEClightgray}{HTML}{F1F5F9}
\definecolor{HEClightblue}{HTML}{E8EDF3}

% --- Packages ---------------------------------------------------------------
\usepackage{amsmath, graphicx, tikz, booktabs, tabularx, hyperref}
\usepackage[inline]{enumitem}
\usepackage{etoolbox}
\usepackage{comment}
\usepackage{needspace}
\usepackage{tcolorbox}
\tcbuselibrary{skins, breakable}
\usepackage{titlesec}
\usepackage{fancyhdr}
\usepackage{lastpage}

% --- Paths ------------------------------------------------------------------
\graphicspath{{../Slides/Figures/}}

% --- Hyperlinks -------------------------------------------------------------
\hypersetup{
   colorlinks,
   linkcolor=HECnavy,
   urlcolor=HECnavy,
   citecolor=HECtext
}

% --- Answer toggle ----------------------------------------------------------
\newbool{answers}
\ifdefined\showanswers
   \booltrue{answers}
\else
   \boolfalse{answers}
\fi

% Answer environment: shown as green box when answers=true, excluded otherwise
\ifbool{answers}{%
   \newtcolorbox{reponse}[1][Réponse]{%
      colback=HECgreen!5,
      colframe=HECgreen,
      boxrule=0pt,
      leftrule=3pt,
      arc=4pt,
      title={\textcolor{HECgreen}{\bfseries #1}},
      coltitle=HECgreen,
      fonttitle=\bfseries,
      attach title to upper={\\[4pt]},
      left=8pt, right=8pt, top=4pt, bottom=4pt,
      before={\par\nopagebreak\medskip\nopagebreak}
   }%
}{%
   \excludecomment{reponse}%
}

% --- Custom box environments (adapted from slides) --------------------------

% Key Insight (green accent)
\newtcolorbox{keyinsight}[1][Point clé]{%
   colback=HECgreen!5,
   colframe=HECgreen,
   boxrule=0pt,
   leftrule=3pt,
   arc=4pt,
   title={\textcolor{HECgreen}{\bfseries #1}},
   coltitle=HECgreen,
   fonttitle=\bfseries,
   attach title to upper={\\[4pt]},
   left=8pt, right=8pt, top=4pt, bottom=4pt,
   breakable
}

% Business Implication (navy accent)
\newtcolorbox{bizimplication}[1][Implication d'affaires]{%
   colback=HECnavy!5,
   colframe=HECnavy,
   boxrule=0pt,
   leftrule=3pt,
   arc=4pt,
   title={\textcolor{HECnavy}{\bfseries #1}},
   coltitle=HECnavy,
   fonttitle=\bfseries,
   attach title to upper={\\[4pt]},
   left=8pt, right=8pt, top=4pt, bottom=4pt,
   breakable
}

% Warning / Important (coral accent)
\newtcolorbox{warning}[1][Important]{%
   colback=HECcoral!5,
   colframe=HECcoral,
   boxrule=0pt,
   leftrule=3pt,
   arc=4pt,
   title={\textcolor{HECcoral}{\bfseries #1}},
   coltitle=HECcoral,
   fonttitle=\bfseries,
   attach title to upper={\\[4pt]},
   left=8pt, right=8pt, top=4pt, bottom=4pt,
   breakable
}

% Definition (gray background)
\newtcolorbox{defbox}[1][Définition]{%
   colback=HEClightgray,
   colframe=HEClightgray,
   boxrule=0pt,
   leftrule=3pt,
   arc=4pt,
   left=8pt, right=8pt, top=4pt, bottom=4pt,
   title={\textcolor{HECnavy}{\bfseries #1}},
   coltitle=HECnavy,
   fonttitle=\bfseries,
   attach title to upper={\\[4pt]},
   breakable
}

% Exercise (coral accent)
\newtcolorbox{exercisebox}[1][Exercice]{%
   colback=HECcoral!5,
   colframe=HECcoral,
   boxrule=0pt,
   leftrule=3pt,
   arc=4pt,
   title={\textcolor{HECcoral}{\bfseries #1}},
   coltitle=HECcoral,
   fonttitle=\bfseries,
   attach title to upper={\\[4pt]},
   left=8pt, right=8pt, top=4pt, bottom=4pt,
   breakable
}

% --- Section formatting (titlesec) ------------------------------------------
\titleformat{\section}
   {\clearpage\Large\bfseries\color{HECnavy}}
   {\thesection.}{0.5em}{}
   [\vspace{2pt}{\color{HECgreen}\rule{2.5cm}{1.5pt}}]
\titlespacing*{\section}{0pt}{1.5em}{0.8em}

\titleformat{\subsection}
   {\large\bfseries\color{HECnavy}}
   {\thesubsection.}{0.5em}{}
\titlespacing*{\subsection}{0pt}{1em}{0.5em}

% --- Header/Footer (fancyhdr) ----------------------------------------------
\pagestyle{fancy}
\fancyhf{}
\renewcommand{\headrulewidth}{0pt}
\renewcommand{\footrulewidth}{0.4pt}
\renewcommand{\footrule}{\hfill{\color{HECgray}\rule{0.9\textwidth}{0.4pt}}\hfill\vspace{4pt}}
\fancyfoot[L]{\small\textcolor{HECgray}{ECON 50803\enspace\textbar\enspace Environnement macroéconomique}}
\fancyfoot[R]{\small\textcolor{HECgray}{\thepage\enspace/\enspace\pageref{LastPage}}}

% --- Enumitem defaults ------------------------------------------------------
\setlist[itemize]{leftmargin=1.5em, itemsep=3pt}
\setlist[enumerate]{leftmargin=1.5em, itemsep=3pt}

% --- MC question helpers ----------------------------------------------------
\newcounter{qcmcounter}
\newcommand{\qcm}[1]{%
   \stepcounter{qcmcounter}%
   \par\vspace{1.2em}%
   \needspace{12\baselineskip}%
   \noindent\textbf{\textcolor{HECnavy}{\theqcmcounter.}}\enspace #1%
   \vspace{0.3em}%
   \nopagebreak[4]%
}

\newcommand{\choix}[1]{%
   \nopagebreak[4]%
   \begin{enumerate}[label=\textcolor{HECgray}{(\alph*)}, leftmargin=2em, itemsep=1pt, topsep=2pt]
      #1
   \end{enumerate}
}

% --- Short answer counter ---------------------------------------------------
\newcounter{shortcounter}
\newcommand{\shortq}[1]{%
   \stepcounter{shortcounter}%
   \par\vspace{1.5em}%
   \needspace{4\baselineskip}%
   \noindent\textbf{\textcolor{HECnavy}{Question \theshortcounter.}}\enspace #1%
   \vspace{0.3em}%
}

% --- VFI (Vrai/Faux/Incertain) helper --------------------------------------
\newcommand{\vfi}[1]{%
   \stepcounter{shortcounter}%
   \par\vspace{1.5em}%
   \needspace{10\baselineskip}%
   \noindent\textbf{\textcolor{HECnavy}{Question \theshortcounter.}}\enspace\textit{Vrai, faux ou incertain?} #1%
   \vspace{0.3em}%
}

% --- Title page command -----------------------------------------------------
\newcommand{\exercisetitlepage}[2]{%
   % #1 = session number, #2 = session title
   \thispagestyle{empty}
   \begin{tikzpicture}[remember picture, overlay]
      % Navy sidebar
      \fill[HECnavy] (current page.north west) rectangle
         ([xshift=0.8cm]current page.south west);
      % Green accent line
      \fill[HECgreen] ([xshift=0.8cm]current page.north west) rectangle
         ([xshift=0.85cm]current page.south west);
   \end{tikzpicture}
   \vspace*{2cm}
   \hspace{0.8cm}%
   \begin{minipage}{0.85\textwidth}
      {\fontsize{28}{34}\selectfont\bfseries\textcolor{HECnavy}{Exercices --- Séance #1}}
      \\[14pt]
      {\Large\textcolor{HECnavy}{#2}}
      \\[14pt]
      {\color{HECgreen}\rule{3cm}{2pt}}
      \\[14pt]
      {\large\textcolor{HECgray}{ECON 50803 --- Environnement macroéconomique}}
      \\[8pt]
      {\large\textcolor{HECgray}{Jean-Félix Brouillette}}
      \\[20pt]
      \ifbool{answers}{%
         {\Large\bfseries\textcolor{HECgreen}{AVEC SOLUTIONS}}%
      }{%
         {\Large\bfseries\textcolor{HECgray}{SANS SOLUTIONS}}%
      }
   \end{minipage}
   \vfill
   \hspace{0.8cm}%
   {\small\textcolor{HECgray}{HEC Montr\'eal\enspace\textbar\enspace MBA}}
   \vspace{1cm}
   \newpage
}

% --- Mathematical commands (from slides) ------------------------------------
\newcommand{\pctchg}[1]{\%\Delta #1}


\begin{document}

\exercisetitlepage{6}{Commerce international et économie ouverte}

% ============================================================================
\section{Questions à choix multiples}
% ============================================================================

\setcounter{qcmcounter}{0}

\qcm{Le pays A a un ratio $(X+M)/Y$ de 70\,\%, alors que le pays B a un ratio de 25\,\%. Toutes choses égales par ailleurs, lequel des deux pays est le plus exposé à un choc commercial mondial\,?}
\choix{
   \item Le pays A, car une plus grande part de son activité est liée au commerce international
   \item Le pays B, car un faible ratio d'ouverture signifie une plus grande fragilité
   \item Les deux pays sont exposés de la même façon, car seul le compte courant compte
   \item On ne peut rien dire sans connaître le taux de change
}

\begin{reponse}
\textbf{(a)} Le pays A est plus exposé, car une plus grande part de sa production et de sa demande est liée aux échanges internationaux. Le ratio $(X+M)/Y$ mesure le \textbf{degré d'ouverture commerciale}. Plus il est élevé, plus un choc sur le commerce mondial, les tarifs ou les chaînes d'approvisionnement risque d'affecter fortement l'économie. Le compte courant mesure un \textbf{solde}; ici, ce qui nous intéresse est plutôt le \textbf{volume} de commerce.
\end{reponse}

\qcm{L'identité fondamentale $CA = S - I$ nous enseigne que le compte courant est déterminé par :}
\choix{
   \item La compétitivité des entreprises nationales sur les marchés mondiaux
   \item L'écart entre l'épargne nationale et l'investissement domestique
   \item Le niveau des droits de douane imposés par les partenaires commerciaux
   \item La valeur du taux de change
}

\begin{reponse}
\textbf{(b)} Le compte courant reflète l'écart entre l'épargne nationale ($S = S_P + S_G$) et l'investissement domestique ($I$). Si $S > I$, le pays prête au reste du monde ($CA > 0$); si $S < I$, il emprunte ($CA < 0$). Cette identité montre que le déficit commercial n'est pas une question de compétitivité ou de tarifs, mais un phénomène \textbf{macroéconomique} lié aux décisions d'épargne et d'investissement.
\end{reponse}

\qcm{Lequel des énoncés suivants est \textbf{vrai} pour une économie fermée?}
\choix{
   \item Le compte courant est positif
   \item $S = I$ à l'équilibre
   \item Le taux d'intérêt est déterminé par le marché mondial
   \item Les exportations nettes sont positives
}

\begin{reponse}
\textbf{(b)} Dans une économie fermée, il n'y a pas de commerce international ($NX = 0$, $CA = 0$). Toute l'épargne nationale finance l'investissement domestique : $S = I$. Le taux d'intérêt d'équilibre ($r^*$) est déterminé par l'intersection des courbes $S(r)$ et $I(r)$ domestiques, \textbf{sans} influence du reste du monde. (a) et (d) sont faux car $NX = CA = 0$ (ni positif ni négatif). (c) décrit une économie ouverte.
\end{reponse}

\qcm{Dans une \textbf{petite économie ouverte}, une hausse des dépenses gouvernementales entraîne :}
\choix{
   \item Une hausse du taux d'intérêt mondial
   \item Une détérioration du compte courant (CA diminue) sans changement du taux d'intérêt
   \item Une hausse du PIB proportionnelle à la hausse de $G$
   \item Aucun effet sur l'économie
}

\begin{reponse}
\textbf{(b)} Dans une petite économie ouverte, le taux d'intérêt est déterminé par le marché mondial ($r^W$) et ne change pas. La hausse de $G$ réduit l'épargne publique ($S_G = T - G$ diminue) $\to$ l'épargne nationale $S$ diminue $\to$ à $r^W$ inchangé, $S < I$ $\to$ le compte courant se détériore ($CA = S - I$ diminue). C'est le mécanisme des \guillemotleft~déficits jumeaux~\guillemotright\ : le déficit budgétaire s'accompagne d'un déficit du compte courant.
\end{reponse}

\qcm{Avant la crise de la zone euro, l'Espagne affichait un fort déficit courant malgré un déficit budgétaire parfois plus modéré qu'on ne le croyait. Quelle explication est la plus cohérente avec l'identité $CA = S - I$\,?}
\choix{
   \item Un boom immobilier et du crédit a fait augmenter fortement l'investissement privé, ce qui a creusé le déficit courant
   \item Les exportations espagnoles ont automatiquement fait baisser l'épargne nationale
   \item Une hausse de l'épargne privée a mécaniquement détérioré le compte courant
   \item L'Espagne a pu réduire son déficit courant en augmentant simplement ses importations
}

\begin{reponse}
\textbf{(a)} Par l'identité $CA = S - I$, un déficit courant peut venir d'une \textbf{hausse de l'investissement} autant que d'une baisse de l'épargne. Dans le cas espagnol, un boom immobilier et du crédit a poussé $I$ vers le haut, ce qui a creusé le déficit courant même sans invoquer d'abord un énorme déficit budgétaire. Les autres réponses violent directement la logique de l'identité ou confondent exportations, importations et épargne.
\end{reponse}

\qcm{L'épargne nationale se décompose en $S = S_P + S_G$ où $S_P = Y - T - C$ et $S_G = T - G$. Une baisse d'impôts ($T\downarrow$) sans changement de $G$ :}
\choix{
   \item Augmente $S_G$ et $S_P$ dans les mêmes proportions
   \item Réduit $S_G$, mais peut augmenter $S_P$ partiellement si les ménages épargnent une partie du gain
   \item N'a aucun effet sur $S$ car les effets sur $S_P$ et $S_G$ s'annulent exactement
   \item Augmente toujours l'épargne nationale
}

\begin{reponse}
\textbf{(b)} Une baisse de $T$ réduit directement $S_G = T - G$ (le déficit budgétaire se creuse). Quant à $S_P = Y - T - C$, le revenu disponible ($Y - T$) augmente, mais les ménages consomment une partie de ce gain ($C\uparrow$). Si les ménages épargnent une fraction de la baisse d'impôts, $S_P$ augmente partiellement. L'effet net sur $S = S_P + S_G$ est généralement \textbf{négatif} : la hausse de $S_P$ ne compense pas la baisse de $S_G$. C'est pourquoi les \guillemotleft~déficits jumeaux~\guillemotright\ sont probables : $S_G\downarrow \to S\downarrow \to CA\downarrow$.
\end{reponse}

\qcm{On observe simultanément que le déficit courant américain se creuse ($CA \downarrow$) alors que le taux d'intérêt réel mondial baisse ($r^W \downarrow$). Quelle explication est la plus cohérente avec ces deux faits\,?}
\choix{
   \item Une hausse des déficits budgétaires américains ($S_G \downarrow$)
   \item Un excès d'épargne dans le reste du monde, qui a fait baisser $r^W$
   \item Un boom d'investissement aux États-Unis ($I \uparrow$)
   \item Une baisse de la productivité américaine
}

\begin{reponse}
\textbf{(b)} Une hausse de l'épargne dans le reste du monde déplace l'offre mondiale de fonds vers la droite $\to$ $r^W \downarrow$. À ce taux mondial plus faible, les États-Unis épargnent relativement moins et investissent davantage, ce qui creuse leur déficit courant ($CA = S - I \downarrow$). Les hypothèses domestiques comme un déficit budgétaire américain plus élevé ou un boom d'investissement peuvent aussi faire baisser $CA$, mais elles tendraient plutôt à faire \textbf{monter} $r^W$, pas à le faire baisser.
\end{reponse}

\qcm{Les États-Unis augmentent fortement leurs dépenses publiques sans hausse équivalente des impôts. Dans le modèle $S$-$I$ d'une \textbf{grande} économie ouverte, quel enchaînement est le plus plausible\,?}
\choix{
   \item L'épargne mondiale se déplace vers la gauche, $r^W$ augmente, et le compte courant américain se détériore
   \item L'épargne mondiale se déplace vers la droite, $r^W$ baisse, et le compte courant américain s'améliore
   \item Le taux d'intérêt mondial ne bouge pas, car le commerce extérieur compense entièrement la baisse d'épargne publique
   \item Les tarifs deviennent plus efficaces pour réduire le déficit courant
}

\begin{reponse}
\textbf{(a)} Une hausse de $G$ non financée par une hausse de $T$ réduit l'épargne publique américaine. Comme les États-Unis sont une \textbf{grande} économie ouverte, cette baisse de l'épargne ne reste pas purement domestique : elle réduit aussi l'offre mondiale de fonds prêtables, ce qui pousse $r^W$ vers le \textbf{haut}. Le compte courant américain tend alors à se \textbf{détériorer} : c'est la logique des déficits jumeaux dans une grande économie ouverte.
\end{reponse}

% ============================================================================
\section{Vrai, faux ou incertain}
% ============================================================================

\setcounter{shortcounter}{0}

\vfi{Un pays qui a un excédent du compte courant est nécessairement en bonne santé économique.}

\begin{reponse}
\textbf{Faux.} Un excédent courant ($CA > 0$) signifie que $S > I$ : le pays épargne plus qu'il n'investit et prête au reste du monde. Mais cela peut refléter un \textbf{manque d'investissement} domestique plutôt qu'une épargne vertueuse. Le Japon, par exemple, affiche des surplus chroniques en partie parce que ses entreprises n'investissent pas suffisamment dans l'économie nationale (faible croissance, vieillissement). L'Allemagne a été critiquée pour des surplus qui reflètent un sous-investissement en infrastructures publiques. Un surplus n'est donc pas automatiquement un signe de santé.
\end{reponse}

\vfi{Un déficit du compte courant signifie qu'un pays vit au-dessus de ses moyens.}

\begin{reponse}
\textbf{Incertain.} Un déficit courant ($CA < 0$) signifie que $S < I$ : le pays investit plus qu'il n'épargne et emprunte au reste du monde. Mais la qualité de cet emprunt dépend de sa \textbf{source} :
\begin{itemize}
   \item \textbf{Déficit \guillemotleft~bon~\guillemotright} : causé par une hausse de l'investissement ($I\uparrow$) financé par des capitaux étrangers. Le pays emprunte pour investir dans des projets \textbf{productifs} qui généreront des revenus futurs. Exemple : les É.-U.\ dans les années 1990 (boom technologique).
   \item \textbf{Déficit \guillemotleft~mauvais~\guillemotright} : causé par une baisse de l'épargne ($S\downarrow$) due à la surconsommation ou au déficit budgétaire. Exemple : la Grèce avant 2010.
\end{itemize}
Le signe du CA ne suffit pas pour juger : c'est la \textbf{cause} du déséquilibre qui détermine s'il est sain ou dangereux.
\end{reponse}

\vfi{Avant la crise de la zone euro, un pays comme l'Espagne pouvait afficher un fort déficit courant même sans grand déficit budgétaire, simplement parce qu'un boom immobilier faisait monter fortement l'investissement privé.}

\begin{reponse}
\textbf{Vrai.} Par l'identité $CA = S - I$, un déficit courant peut apparaître même si l'épargne publique n'est pas particulièrement faible. Il suffit que l'\textbf{investissement domestique} augmente fortement relativement à l'épargne nationale. C'est précisément ce qui s'est produit dans certains pays de la périphérie de la zone euro avant 2008 : un boom immobilier et du crédit a fait monter $I$, ce qui a creusé le déficit courant sans qu'il soit nécessaire d'invoquer d'abord un grand déficit budgétaire.
\end{reponse}

\vfi{Si le Canada impose des tarifs sur les importations américaines, le taux d'intérêt mondial ($r^W$) augmentera.}

\begin{reponse}
\textbf{Faux.} Le Canada est une \textbf{petite économie ouverte} : ses décisions n'affectent pas $r^W$. Les tarifs canadiens changeraient la composition du commerce canadien, mais pas les conditions financières mondiales. Le taux $r^W$ est déterminé par l'offre et la demande mondiales de fonds prêtables, qui dépendent principalement des décisions d'épargne et d'investissement des \textbf{grandes} économies (É.-U., Chine, zone euro).
\end{reponse}

\vfi{Lorsqu'un pays s'ouvre au commerce international et que $r^W < r^*$ (taux d'intérêt d'autarcie), le pays devient prêteur net sur les marchés internationaux.}

\begin{reponse}
\textbf{Faux.} C'est l'inverse. Si $r^W < r^*$, le capital étranger est \textbf{bon marché} par rapport au rendement que les entreprises domestiques peuvent obtenir. Les entreprises empruntent davantage à l'étranger pour investir $\to$ $I > S$ au taux $r^W$ $\to$ $CA = S - I < 0$ $\to$ le pays est \textbf{emprunteur net}. Pour qu'un pays soit prêteur net, il faut $r^W > r^*$ : le rendement mondial est plus attractif que le rendement domestique, et les épargnants domestiques préfèrent prêter à l'étranger.
\end{reponse}

\vfi{L'expansion rapide des dépenses publiques aux É.-U.\ et en Europe pendant et après la pandémie devrait nous inquiéter du risque de détérioration des balances commerciales de ces économies, toutes choses égales par ailleurs.}

\begin{reponse}
\textbf{Vrai.} Dans le modèle $S$-$I$, une hausse de $G$ (sans hausse de $T$) réduit l'épargne publique ($S_G = T - G \downarrow$), ce qui réduit l'épargne nationale ($S\downarrow$). À $I$ inchangé, $CA = S - I$ se détériore. Pour une grande économie comme les É.-U., le déplacement de $S$ vers la gauche fait aussi monter $r^W$, ce qui affecte le reste du monde. C'est le mécanisme des \guillemotleft~déficits jumeaux~\guillemotright\ vu en cours : déficit budgétaire $\to$ déficit courant. Le compte courant américain est effectivement passé d'environ $-2\,\%$ du PIB avant la pandémie à $-3.5\,\%$ en 2022.
\end{reponse}

\vfi{Si les Américains augmentaient leur taux d'épargne, le déficit commercial américain se réduirait.}

\begin{reponse}
\textbf{Vrai.} Par l'identité $CA = S - I$, une hausse de $S$ (à $I$ constant) augmente le CA et réduit le déficit. C'est exactement le message \guillemotleft~Make America Save Again~\guillemotright\ vu en cours : le déficit commercial reflète un manque d'épargne, pas un excès d'importations. Pour réduire le déficit, il faut épargner davantage --- pas imposer des tarifs. Cela peut se faire via une hausse de l'épargne privée (ménages) ou publique (réduction du déficit budgétaire).
\end{reponse}

% ============================================================================
\section{Questions à développement}
% ============================================================================

\setcounter{shortcounter}{0}

\shortq{En 2022, la hausse du prix du pétrole améliore fortement les revenus d'exportation de deux pays producteurs. Le pays A dépense rapidement cette manne par une hausse de la consommation publique et privée. Le pays B place plutôt une grande partie de ces revenus dans un fonds souverain, sans augmenter beaucoup sa consommation ni son investissement domestique à court terme. Utilisez l'identité $CA = S - I$ pour comparer les deux pays.}

\begin{enumerate}[label=(\alph*)]
   \item Dans quel pays l'épargne nationale augmente le plus\,? Expliquez.
   \item En supposant que l'investissement domestique change peu à court terme, quel pays verra son compte courant s'améliorer le plus\,? Pourquoi\,?
   \item Quel pays ressemble davantage à un pays comme la Norvège dans sa gestion d'un choc positif sur les prix de l'énergie\,? Pourquoi\,?
   \item Supposons maintenant que le pays A utilise une partie de cette manne pour financer progressivement des infrastructures productives domestiques. Comment cela modifie-t-il votre diagnostic sur $I$ et sur le compte courant\,?
\end{enumerate}

\begin{reponse}
\textbf{(a)} Le \textbf{pays B} voit son épargne nationale augmenter le plus. Comme il ne transforme pas immédiatement la manne pétrolière en consommation ou en dépenses domestiques, une plus grande part du revenu supplémentaire est \textbf{épargnée}. Dans le pays A, au contraire, une large part du gain est rapidement consommée par le secteur privé ou public, donc la hausse de $S$ est plus faible.

\textbf{(b)} À investissement domestique à peu près inchangé à court terme, le pays dont l'épargne augmente le plus voit aussi son \textbf{compte courant} s'améliorer le plus, car
\[
CA = S - I.
\]
Le \textbf{pays B} verra donc son compte courant s'améliorer davantage. Le pays A peut aussi voir son compte courant s'améliorer grâce à la hausse des revenus d'exportation, mais beaucoup moins si cette manne est immédiatement dépensée.

\textbf{(c)} Le \textbf{pays B} ressemble davantage à la \textbf{Norvège}. L'idée centrale est qu'une partie importante des revenus tirés des ressources naturelles est mise de côté dans un fonds souverain plutôt que transformée immédiatement en demande intérieure. Macroéconomiquement, cela revient à augmenter fortement l'épargne nationale, donc à améliorer le compte courant.

\textbf{(d)} Si le pays A finance progressivement des infrastructures productives, l'\textbf{investissement domestique} ($I$) augmente. Cela change le diagnostic : même si le revenu pétrolier augmente aussi l'épargne, une partie de cette amélioration est maintenant absorbée par une hausse de $I$. Le compte courant peut encore s'améliorer, mais \textbf{moins} que si les revenus étaient simplement épargnés. En d'autres mots, la manne est davantage transformée en capital domestique qu'en créances sur le reste du monde.
\end{reponse}

\shortq{En utilisant le diagramme $S$-$I$ en économie ouverte (deux panneaux : Monde et Domestique), analysez l'effet d'un vieillissement accéléré de la population dans l'économie domestique. Distinguez les cas (i) d'une petite économie ouverte et (ii) d'une grande économie ouverte.}

\begin{reponse}
\textbf{Mécanisme commun :} Le vieillissement réduit l'épargne nationale. Les retraités \textbf{désépargnent} (ils consomment leur épargne accumulée), et la proportion de retraités dans la population augmente. La courbe $S$ du pays se déplace vers la \textbf{gauche}.

\textbf{(i) Petite économie ouverte :}
\begin{itemize}
   \item Le pays ne peut pas affecter $r^W$ $\to$ $r^W$ reste \textbf{inchangé}
   \item Au même $r^W$, $S$ diminue mais $I$ reste stable
   \item $CA = S - I$ se \textbf{détériore} (le déficit courant augmente ou le surplus diminue)
   \item Le pays emprunte davantage au reste du monde pour financer l'écart $S < I$
\end{itemize}

\textbf{(ii) Grande économie ouverte :}
\begin{itemize}
   \item La baisse de $S$ domestique réduit aussi $S^W$ (car le pays est assez grand pour peser sur l'épargne mondiale)
   \item Dans le panneau Monde : $S^W$ se déplace vers la gauche $\to$ $r^W$ \textbf{augmente}
   \item Dans le panneau Domestique : au $r^W$ plus élevé, $S$ diminue \textbf{moins} (incitation à épargner plus) et $I$ diminue un peu (coût du capital plus élevé)
   \item $CA$ se détériore, mais \textbf{moins} qu'en petite économie (la hausse de $r^W$ modère l'effet)
   \item Le reste du monde voit $r^W$ augmenter $\to$ son CA \textbf{s'améliore}
\end{itemize}

\textbf{Implication d'affaires :} Le vieillissement démographique dans les économies avancées (Japon, Europe, Canada) pousse les taux d'intérêt à la hausse et aggrave les déficits extérieurs. C'est un facteur structurel de long terme, pas un choc temporaire.
\end{reponse}

\end{document}
